% 来源：1020-12.pdf，第 40 页。本次只整理定义与两组基底的变换。
\section{Clebsch--Gordan 系数}
固定 $j_1,j_2$，无耦合基 $\left\{\Ket{j_1,m_1;j_2,m_2}\right\}$ 与耦合基
$\left\{\Ket{j_1,j_2;j,m}\right\}$ 都是同一个
$\left(2j_1+1\right)\left(2j_2+1\right)$ 维空间中的基底。插入耦合基的完备关系，得到
\begin{align}
\Ket{j_1,m_1;j_2,m_2}
&=\hat I\Ket{j_1,m_1;j_2,m_2}\notag\\
&=\sum_{j,m}\Ket{j_1,j_2;j,m}
\Braket{j_1,j_2;j,m|j_1,m_1;j_2,m_2}.
\end{align}
反过来，插入无耦合基的完备关系，
\begin{align}
\Ket{j_1,j_2;j,m}
&=\hat I\Ket{j_1,j_2;j,m}\notag\\
&=\sum_{m_1,m_2}\Ket{j_1,m_1;j_2,m_2}
\Braket{j_1,m_1;j_2,m_2|j_1,j_2;j,m}.
\end{align}
两组基矢之间的重叠称为 Clebsch--Gordan 系数（简称 C--G 系数）。为固定记号，定义
\begin{equation}
C_{j_1m_1,j_2m_2}^{jm}
=\Braket{j_1,m_1;j_2,m_2|j_1,j_2;j,m}.
\end{equation}
反向变换中的重叠是其复共轭：
\begin{equation}
\Braket{j_1,j_2;j,m|j_1,m_1;j_2,m_2}
=\left(C_{j_1m_1,j_2m_2}^{jm}\right)^*.
\end{equation}
在通常的实系数相位约定下，两者相同；未选定相位约定前应保留共轭。照片本次内容到此为止。
